\subsection{Conclusions and Future Directions}

This study compared classical machine learning ensembles with a graph neural network on a small TgDHFR data set ($N = \NCompounds$), under a leakage-free nested cross-validation protocol, and applied the best classical model to a library of FDA-approved drugs with an orthogonal docking check. Three conclusions follow.

First, on this data set a classical stacking ensemble on 2D Morgan fingerprints performed at least as well as the alternatives we tested. It reached $R^2 = \MorganRTwoMean \pm \MorganRTwoStd$, against $\ChemPropDMPNNFoldRTwoMean \pm \ChemPropDMPNNFoldRTwoStd$ for a default D-MPNN and $\GnnTunedRTwoMean \pm \GnnTunedRTwoStd$ for one tuned in a nested search; the difference ($\Delta R^2 = \GnnTunedVsMorganDiffRTwo$ for the tuned network) was small but survived multiple-comparison correction, and equivalence with the Morgan ensemble could not be established in either metric, so the network as usually trained is behind the fingerprint-based ensembles rather than interchangeable with them. The gap is one of training budget and not of architecture: averaging five networks narrows it, and a D-MPNN initialized from a pretrained model reaches $\ChemeleonRTwoMean$ and is the only model equivalent to the Morgan ensemble in both metrics, at \ChemeleonCostFactor\ times the computational cost and without exceeding it. Adding 3D descriptors to the fingerprint-based sets changed nothing detectable. Feature selection and Gaussian-noise augmentation, which the original study credited with large gains, did not improve a classical learner under the same hold-out protocol. Those gains are accounted for by information leaks: removing them one at a time from the published code takes its $R^2$ from \KarLadderRungZeroRTwoMean\ to \KarLadderRungFourRTwoMean, and the same selection applied leakily to a classical learner reproduces them (\ExpBleakyRTwoMean\ against \ExpBRTwoMean). These findings are limited to a single target and a comparison in which the network search was modest; they do not show that graph networks cannot match fixed fingerprints. The classical protocol was also not cheaper to run: the benchmark took longer than the network's, and its advantage lies in accuracy that is not lower, simpler inputs, and a deployed model that needs only a subset of fingerprint bits.

Second, nested cross-validation and a careful curation of the data (removal of duplicates and of two extreme potency values, with a sensitivity analysis of the latter) were necessary to obtain estimates that can be compared across models. Our random-split estimates are also optimistic: with scaffold-grouped splits every model lost between \ScaffoldChemeleonDropRTwo\ and \ScaffoldRdkitTwoDDropRTwo\ units of $R^2$ (to $\ScaffoldMorganRTwoMean$ for the Morgan ensemble and $\ScaffoldGnnRTwoMean$ for the D-MPNN), and the data could not resolve the differences between them. The loss was smallest for the pretrained network, so what representation learning offers on this data set is robustness to a change of chemotype rather than accuracy within one.

Third, combining ML prediction with docking produced a short, reproducible list of hypotheses and not validated leads. Only \ScreenNDomain{} of the \ScreenNStructural{} scored compounds fell inside the model's applicability domain, the agreement between predicted potency and docking energy was weak ($\rho = \DockSpearmanRho$, $p = \DockSpearmanP$), and the compounds leading the ranking by potency are antifolates that also inhibit the human enzyme, whose selectivity was not assessed.

Future work should therefore (i) dock the consensus hits against human DHFR and test them in enzymatic assays against TgDHFR, (ii) evaluate temporal splits, (iii) compare against larger searches and pretrained graph models, and (iv) apply the pipeline to other targets before claiming that it generalizes.
